%% kapa: exam=1 round=30 year=2019 period=1 form=A subject=civil_law subject_ko=민법 seq=38 number=38
%% source: Q-Net 감정평가사 기출문제 (cst00309) artlSeq=5205205 file=extracted/1008/5205205/2019년 제30회 감정평가사 1차시험 문제지/2019년 제30회 감정평가사 1차시험 문제지/2019년도 제30회 감정평가사 1차 1교시 A형.hwp
%% conversion: kordoc 4.24.51 -> tools/gen1.py (auto); review=proofread-sonnet; flags=none
\begin{question}{38}
권리질권에 관한 설명으로 옳은 것은?
\begin{choices}
  \item 부동산의 사용을 목적으로 하는 권리도 질권의 목적이 될 수 있다.
  \item 질권자는 질권의 목적이 된 채권을 직접 청구할 수 없다.
  \item 지명채권을 목적으로 한 질권은 제3채무자에게 질권설정의 사실을 통지하여야 성립할 수 있다.
  \item 입질된 채권의 목적물이 금전 이외의 물건인 때에는 질권자는 그 변제를 받은 물건에 대하여 질권을 행사할 수 있다.
  \item 지시채권을 목적으로 한 질권의 설정은 배서없이 증서를 교부하더라도 그 효력이 생긴다.
\end{choices}
\end{question}
\answer{4}
